\scanpage{213}
\mspage{1 bis}
% archive sheet number: 212

\deleted{$\exists$ fonction}

\textcircled{3} On peut trouver, pour tt $n\in\mathbb N$, la classe des ``morphismes de degré $n$'', de façon que

\uline{a)} Stabilité par changt de base

\uline{b)} additivité finie

\uline{c)} Un isom. est de degré 1

\uline{f)} Si $f:X\longrightarrow Y$, alors $Y=\coprod_{n\in\mathbb N}Y_n$, somme directe disjointe, effective, universelle, tels que $X_n=X\times_Y Y_n\longrightarrow Y_n$ soit de degré $n$.

\uline{d)} Si $f:X\longrightarrow Y$ est de degré $n\ge1$, alors $f$ est un morphisme de descente effective [donc de descente effective \uline{universelle} grâce à a)]

\uline{e)} Si $f:X\longrightarrow Y$ est de degré 0, alors $X=\emptyset_C$

\textcircled{2} Tout monomorphisme est un isomorphisme avec un sommand direct universel

\textcircled{1} Existence de $\underleftarrow{\mathrm{lim}}$ finies

\deleted{\textcircled{\unsure{1|0}}} \unsure{$e$|\deleted{$e$}}

\textcircled{0} $C$ est une $\mathcal U$-catégorie (i.e. les $\mathrm{Hom}(X,Y)$ sont $\mathcal U$-petits) et l'ens des classes d'isom. d'objets de $C$ est $\mathcal U$-petit.

\unsure{'} Implique condition A(iii). En effet, à prouver que pour $X$ donné, les $S\longrightarrow e$ tels que $X_S$ constant fini sur $S$ font famille \unsure{de|de la} descente eff. universelle. Grâce à 3 f), on est ramené à prouver l'analogue pour $X\longrightarrow Y$ au lieu de $X\longrightarrow e$, mais avec $X\longrightarrow Y$ de degré $n$. \inserted{à trouver un morphisme de \unsure{desc. eff. univ.|descente eff. universelle} \deleted{\unsure{ou}} qui trivialise $X$} \inserted{Si} $n=0$, on conclut par 3 e que $X=\emptyset$, on prend $Y\longrightarrow Y$ ($X=\unsure{\emptyset_{/Y}|\emptyset_Y}$ !). Si $n>0$, on procède par récurrence, supposant que

